\chapter{Amplitwists: Scaling and Rotation Combined}
\label{ch:ch15-amplitwists}

\section{Scaling and Rotation as a Single Transformation}

In Chapter~1 we introduced scaling transformations of the form
$(x,y)\mapsto(kx,ky)$, which enlarge or shrink figures while preserving
their shape. In Chapter~8 we introduced rotations, represented by the
matrix
\[
R(\theta)=
\begin{pmatrix}
\cos\theta & -\sin\theta\\
\sin\theta & \cos\theta
\end{pmatrix},
\]
which preserve lengths while changing orientation. These transformations
were studied separately: scaling changes size but not direction, and
rotation changes direction but not size. A natural question therefore
arises, whether scaling and rotation can be combined into a single
transformation. The answer is yes, and this chapter develops a unified
operator that performs both actions simultaneously.

\section{Scaling as a Matrix}

Recall that scaling every vector by a factor $s$ can be written as
\[
S(s)=
\begin{pmatrix}
s & 0\\
0 & s
\end{pmatrix}.
\]
Applying this matrix to a vector $\mathbf v = \begin{pmatrix} x\\ y
\end{pmatrix}$ gives $S(s)\mathbf v = \begin{pmatrix} sx\\ sy
\end{pmatrix}$: every coordinate is multiplied by the same factor, so
every length is multiplied by $s$.

\begin{example}
Let $\mathbf v = \begin{pmatrix} 3\\ 4 \end{pmatrix}$. Applying the
scaling matrix $S(2) = \begin{pmatrix} 2&0\\ 0&2 \end{pmatrix}$ produces
$S(2)\mathbf v = \begin{pmatrix} 6\\ 8 \end{pmatrix}$. The original
length was $\sqrt{3^2+4^2}=5$; the new length is $\sqrt{6^2+8^2}=10$. The
vector has doubled in length while preserving its direction.
\end{example}

\section{Rotation Revisited}

A rotation through an angle $\theta$ is represented by $R(\theta)$ as
above; applying $R(\theta)$ changes the direction of a vector without
changing its length. From Chapter~8, recall the composition rule
$R(\theta_1)R(\theta_2) = R(\theta_1+\theta_2)$, which will play a
central role in what follows.

\section{The Amplitwist}

\begin{definition}
An \emph{amplitwist} is a transformation obtained by combining a scaling
factor $s>0$ with a rotation through an angle $\theta$: $A = sR(\theta)$.
\end{definition}

\begin{historicalnote}{The Origin of the Term ``Amplitwist''}
The term \emph{amplitwist} was coined by Tristan Needham in his textbook
\emph{Visual Complex Analysis} \cite{needham1997}, where it names the
local action of a complex-differentiable function: near a point $a$, such
a function acts on an infinitesimally small displacement as multiplication
by a fixed complex number, simultaneously amplifying its length and
twisting its direction, hence \emph{amplitwist}. This chapter borrows
Needham's term and geometric picture but restricts it to a single global
transformation of the plane, $A = sR(\theta)$, applied everywhere at once,
rather than a transformation that varies infinitesimally from point to
point as it does in Needham's original setting. Chapter~17's treatment of
derivatives, and in particular its identification of differentiation with
a quarter-turn rotation, is the closest this book comes to Needham's
full picture: a differentiable function's derivative at a point is, in
exactly this sense, the amplitwist that best approximates the function
near that point.
\end{historicalnote}

Expanding the matrix gives
\[
A=
\begin{pmatrix}
s\cos\theta & -s\sin\theta\\
s\sin\theta & s\cos\theta
\end{pmatrix}.
\]
The parameter $s$ controls size while $\theta$ controls orientation.

\begin{example}
Let $A = 2R(30^\circ)$. Applying $A$ to $\mathbf v = \begin{pmatrix} 1\\
0 \end{pmatrix}$ gives
\[
A\mathbf v=
\begin{pmatrix}
2\cos30^\circ\\
2\sin30^\circ
\end{pmatrix}
=
\begin{pmatrix}
\sqrt3\\
1
\end{pmatrix}.
\]
The vector has been rotated by $30^\circ$ and its length has doubled.
\end{example}

\section{Geometric Meaning}

An amplitwist preserves all angles and shape relationships; it changes
only the overall size of a figure and its orientation, altering neither
the relative proportions between the figure's parts nor the angles
between its edges. When $\theta = 0$, the amplitwist reduces to a pure
scaling transformation, $A = sR(0) = S(s)$, so every similarity
transformation fixing the origin is an amplitwist.

\section{Composition of Amplitwists}

The central result of this chapter is that amplitwists compose in a
remarkably simple way.

\begin{theorem}
If $A_1 = s_1R(\theta_1)$ and $A_2 = s_2R(\theta_2)$, then $A_1A_2 =
s_1s_2R(\theta_1+\theta_2)$.
\end{theorem}

\begin{proof}
Using matrix multiplication, $A_1A_2 = (s_1R(\theta_1))(s_2R(\theta_2))$.
Since scalars commute with matrices, $A_1A_2 = s_1s_2
R(\theta_1)R(\theta_2)$. By the rotation composition theorem of
Chapter~8, $R(\theta_1)R(\theta_2) = R(\theta_1+\theta_2)$, so $A_1A_2 =
s_1s_2R(\theta_1+\theta_2)$: amplitudes multiply while angles add.
\end{proof}

\begin{algorithmproc}{How to Compose Two Amplitwists}
Given amplitwists $(s_1,\theta_1)$ and $(s_2,\theta_2)$, multiply the
amplitudes to obtain $s = s_1s_2$, add the angles to obtain $\theta =
\theta_1+\theta_2$, and write the result as the amplitwist $(s,\theta)$.
\end{algorithmproc}

\begin{practicenow}
Compose the amplitwists $(3,50^\circ)$ and $\left(\dfrac12,
110^\circ\right)$, giving your answer both in amplitwist coordinates and
as $sR(\theta)$.
\end{practicenow}

\section{Identity and Inverses}

The identity amplitwist is $I = 1\cdot R(0)$; applying $I$ leaves every
vector unchanged.

\begin{theorem}
The inverse of $A = sR(\theta)$ is $A^{-1} = \dfrac1s R(-\theta)$.
\end{theorem}

\begin{proof}
Compute $AA^{-1} = sR(\theta)\left(\dfrac1s R(-\theta)\right)$, which
simplifies to $R(\theta)R(-\theta)$. Using rotation composition,
$AA^{-1} = R(0) = I$, so $A^{-1} = \dfrac1s R(-\theta)$.
\end{proof}

\section{Amplitwist Coordinates}

Every amplitwist can be represented by the ordered pair $(s,\theta)$, and
composition becomes $(s_1,\theta_1)\circ(s_2,\theta_2) =
(s_1s_2,\theta_1+\theta_2)$. Two different arithmetic operations appear
simultaneously in this single rule: scales are multiplied while angles
are added. This unusual combination will reappear, in a new guise, in the
next chapter.

\section{Repeated Amplitwists and Spirals}

Repeated application of an amplitwist produces a simple pattern.

\begin{theorem}
For every positive integer $n$, $A^n = s^nR(n\theta)$.
\end{theorem}

\begin{proof}
We proceed by induction. For $n=1$, $A^1 = sR(\theta)$, which is true.
Assume $A^n = s^nR(n\theta)$; then $A^{n+1} = A^nA$, and substituting the
inductive hypothesis gives $A^{n+1} = s^nR(n\theta)\cdot sR(\theta)$.
Using the composition theorem, $A^{n+1} = s^{n+1}R((n+1)\theta)$, so the
formula holds for $n+1$, completing the induction.
\end{proof}

Repeated amplitwists generate spiral trajectories: if $s>1$, the spiral
expands outward; if $0<s<1$, it contracts inward; and if $s=1$, the
motion becomes pure rotation around a circle, since the amplitude never
changes at all.

\section{Applications}

Amplitwists appear naturally in many fields. In computer graphics,
objects are frequently rotated and scaled simultaneously; in robotics,
coordinate systems are continually transformed as machines move through
space; in signal processing, amplitude and phase changes occur together
in oscillatory systems; and in fractal geometry and iterated function
systems, repeated amplitwists generate self-similar structures and
spiral patterns.

\section{Looking Ahead}

The composition law $(s_1,\theta_1)\circ(s_2,\theta_2) =
(s_1s_2,\theta_1+\theta_2)$ may seem unusual at first. In the next
chapter we will discover that complex numbers obey exactly this same
rule: multiplying complex numbers turns out to be nothing more than
amplitwist composition written in algebraic form. Complex numbers
therefore provide a numerical language for the geometry developed in
this chapter.

\section{Exercises}

\subsection*{Basic Computations}
\begin{exercise}
Let $A = 2R(45^\circ)$. Apply $A$ to the vector $\begin{pmatrix} 1\\ 0
\end{pmatrix}$, and describe both the change in length and the change in
direction.
\end{exercise}

\begin{exercise}
Let $A = 3R(90^\circ)$. Apply $A$ to $\begin{pmatrix} 2\\ 1
\end{pmatrix}$.
\end{exercise}

\begin{exercise}
Write the amplitwist $4R(30^\circ)$ as a single $2\times2$ matrix.
\end{exercise}

\begin{exercise}
Write the amplitwist $\dfrac12 R(120^\circ)$ as a single $2\times2$
matrix.
\end{exercise}

\begin{exercise}
Determine the image of $\begin{pmatrix} 4\\ 0 \end{pmatrix}$ under the
amplitwist $2R(60^\circ)$.
\end{exercise}

\subsection*{Composition of Amplitwists}
\begin{exercise}
Compute $(2,30^\circ)\circ(3,40^\circ)$, expressing the answer both in
amplitwist coordinates and as $sR(\theta)$.
\end{exercise}

\begin{exercise}
Compute $(5,15^\circ)\circ\left(\dfrac12,75^\circ\right)$.
\end{exercise}

\begin{exercise}
Verify directly by matrix multiplication that $2R(30^\circ)\cdot
3R(45^\circ) = 6R(75^\circ)$.
\end{exercise}

\begin{exercise}
Compose the amplitwists $(4,120^\circ)$ and $\left(\dfrac14,
-120^\circ\right)$. What transformation results?
\end{exercise}

\begin{exercise}
Find the amplitwist obtained by composing $(2,90^\circ)$,
$(2,90^\circ)$, and $(2,90^\circ)$.
\end{exercise}

\subsection*{Identity and Inverses}
\begin{exercise}
Find the inverse of $A = 5R(40^\circ)$.
\end{exercise}

\begin{exercise}
Find the inverse of $A = \dfrac13 R(-120^\circ)$.
\end{exercise}

\begin{exercise}
Verify directly that $AA^{-1} = I$ for $A = 2R(60^\circ)$.
\end{exercise}

\begin{exercise}
Determine the amplitwist whose composition with $(7,-25^\circ)$ produces
the identity.
\end{exercise}

\begin{exercise}
Explain geometrically why the inverse amplitwist must reverse both the
scaling and the rotation.
\end{exercise}

\subsection*{Amplitwist Coordinates}
\begin{exercise}
Express the composition $(3,20^\circ)\circ(4,10^\circ)$ using amplitwist
coordinates only.
\end{exercise}

\begin{exercise}
Compute $(2,45^\circ)^2$.
\end{exercise}

\begin{exercise}
Compute $(2,45^\circ)^3$.
\end{exercise}

\begin{exercise}
Compute $(2,45^\circ)^4$. What familiar transformation appears?
\end{exercise}

\begin{exercise}
Explain why amplitudes multiply but angles add when amplitwists are
composed.
\end{exercise}

\subsection*{Repeated Amplitwists and Spirals}
\begin{exercise}
Let $A = 1.2R(30^\circ)$. Compute $A$, $A^2$, $A^3$, $A^4$, and describe
the resulting sequence geometrically.
\end{exercise}

\begin{exercise}
Let $A = 0.8R(45^\circ)$. Compute $A^5$. Is the resulting orbit
expanding or contracting?
\end{exercise}

\begin{exercise}
Let $A = R(72^\circ)$. Compute $A^5$. What does this reveal about
repeated rotations?
\end{exercise}

\begin{exercise}
Starting with $\mathbf v = \begin{pmatrix} 1\\ 0 \end{pmatrix}$, apply
$A = 1.1R(40^\circ)$ repeatedly, and sketch the first six points.
\end{exercise}

\begin{exercise}
Describe the qualitative behavior of repeated amplitwists when $s>1$,
$0<s<1$, and $s=1$.
\end{exercise}

\subsection*{Proofs and Theory}
\begin{exercise}
Prove by induction that $A^n = s^nR(n\theta)$ for every positive integer
$n$.
\end{exercise}

\begin{exercise}
Show that the set of all amplitwists is closed under composition.
\end{exercise}

\begin{exercise}
Show that the composition of two amplitwists is always an amplitwist.
\end{exercise}

\begin{exercise}
Show that every amplitwist has a unique inverse.
\end{exercise}

\begin{exercise}
Prove that the identity amplitwist is $(1,0)$.
\end{exercise}

\subsection*{Challenge Problems}
\begin{exercise}
Let $A = sR(\theta)$. Determine all amplitwists satisfying $A^2 = I$.
\end{exercise}

\begin{exercise}
Determine all amplitwists satisfying $A^3 = I$.
\end{exercise}

\begin{exercise}
Suppose $A = sR(\theta)$ and $B = tR(\phi)$. Show that $AB = BA$.
\end{exercise}

\begin{exercise}
Prove that the amplitwist coordinates form a group under composition.
\end{exercise}

\begin{exercise}
Show that every orientation-preserving similarity transformation fixing
the origin can be written uniquely as $A = sR(\theta)$.
\end{exercise}

\begin{exercise}
Explain why amplitwist composition resembles a new arithmetic. What
quantities are being multiplied, and what quantities are being added?
\end{exercise}
